% R3 journal manuscript: editorial and LaTeX-portability refinement.
% Compile with pdfLaTeX (twice for cross-references); IEEEtran supplies the text fonts.
% Compatibility layer for the actual Tectonic/XeTeX build only.
% pdfLaTeX ignores every conditional definition below.
\ifdefined\XeTeXversion
  \newif\ifTKTermes
  \TKTermestrue
  \IfFileExists{texgyretermes-regular.otf}{}{\TKTermesfalse}
  \IfFileExists{texgyretermes-italic.otf}{}{\TKTermesfalse}
  \IfFileExists{texgyretermes-bold.otf}{}{\TKTermesfalse}
  \IfFileExists{texgyretermes-bolditalic.otf}{}{\TKTermesfalse}
  \ifTKTermes
    \def\TKRomanRegular{[texgyretermes-regular.otf]}
    \def\TKRomanItalic{[texgyretermes-italic.otf]}
    \def\TKRomanBold{[texgyretermes-bold.otf]}
    \def\TKRomanBoldItalic{[texgyretermes-bolditalic.otf]}
  \else
    \def\TKRomanRegular{Times New Roman/OT}
    \def\TKRomanItalic{Times New Roman Italic/OT}
    \def\TKRomanBold{Times New Roman Bold/OT}
    \def\TKRomanBoldItalic{Times New Roman Bold Italic/OT}
  \fi
  \DeclareFontFamily{TU}{ptm}{}
  \DeclareFontShape{TU}{ptm}{m}{n}{<->"\TKRomanRegular:mapping=tex-text"}{}
  \DeclareFontShape{TU}{ptm}{m}{it}{<->"\TKRomanItalic:mapping=tex-text"}{}
  \DeclareFontShape{TU}{ptm}{bx}{n}{<->"\TKRomanBold:mapping=tex-text"}{}
  \DeclareFontShape{TU}{ptm}{bx}{it}{<->"\TKRomanBoldItalic:mapping=tex-text"}{}
  \DeclareFontShape{TU}{ptm}{b}{n}{<->ssub*ptm/bx/n}{}
  \DeclareFontShape{TU}{ptm}{b}{it}{<->ssub*ptm/bx/it}{}
  \DeclareFontShape{TU}{ptm}{m}{sl}{<->ssub*ptm/m/it}{}
  \DeclareFontShape{TU}{ptm}{bx}{sl}{<->ssub*ptm/bx/it}{}
  \ifTKTermes
    \DeclareFontShape{TU}{ptm}{m}{sc}{<->"\TKRomanRegular:mapping=tex-text;+smcp"}{}
  \else
    \DeclareFontShape{TU}{ptm}{m}{sc}{<->ssub*ptm/m/n}{}
  \fi
\fi

\documentclass[journal]{IEEEtran}
\usepackage{cite}
\usepackage{amsmath,amssymb,amsfonts}
% Must follow amsmath: Tectonic's cached XeTeX lacks cmss10.
\ifdefined\XeTeXversion
  \DeclareFontFamily{TU}{qhv}{}
  \IfFileExists{texgyreheros-regular.otf}{
    \DeclareFontShape{TU}{qhv}{m}{n}{<->"[texgyreheros-regular.otf]"}{}
  }{
    \DeclareFontShape{TU}{qhv}{m}{n}{<->"Arial"}{}
  }
  \DeclareMathAlphabet{\mathsf}{TU}{qhv}{m}{n}
\fi

\usepackage{graphicx}
\usepackage[utf8]{inputenc}
\usepackage[T1]{fontenc}
\usepackage{microtype,amsthm,algorithm,algorithmic}
\usepackage[hidelinks]{hyperref}
\newtheorem{assumption}{Assumption}
\newtheorem{theorem}{Theorem}
\newtheorem{proposition}{Proposition}
\newtheorem{corollary}{Corollary}
\newtheorem{remark}{Remark}
\interdisplaylinepenalty=2500
\emergencystretch=2em
\newcommand{\vphi}{\boldsymbol{\phi}}
\newcommand{\vp}{\mathbf{p}}
\newcommand{\vb}{\mathbf{b}}
\newcommand{\vx}{\mathbf{x}}
\newcommand{\vu}{\mathbf{u}}
\newcommand{\vs}{\mathbf{s}}
\newcommand{\vz}{\mathbf{z}}
\newcommand{\vc}{\mathbf{c}}
\newcommand{\mH}{\mathbf{H}}
\newcommand{\mA}{\mathbf{A}}
\newcommand{\mB}{\mathbf{B}}
\newcommand{\mU}{\mathbf{U}}
\newcommand{\mV}{\mathbf{V}}
\newcommand{\mW}{\mathbf{W}}
\newcommand{\mX}{\mathbf{X}}
\newcommand{\R}{\mathbb{R}}
\newcommand{\norm}[1]{\left\lVert #1\right\rVert}
\newcommand{\col}{\operatorname{col}}

\begin{document}
\title{Transformer--Koopman: History-Conditioned Scheduled Prediction and Frozen-LTV Model Predictive Control}
\author{Bolin Tao}
\markboth{IEEE Journal Draft}{Tao: Transformer--Koopman}
\maketitle

\begin{abstract}
Finite measurement histories can distinguish dynamical conditions that a single observation cannot, but history-dependent prediction generally couples future dynamics to the candidate control inputs. This paper combines a measurement-preserving temporal Transformer lift with low-rank transition scheduling driven by recursively updated history. A nonredundant augmented state carries the lifted coordinates and the measurement and input caches. One nonlinear nominal rollout then supplies a frozen linear time-varying model for a single convex quadratic program at each sample. A reference-relative objective and a shift of the returned sequence with a reference tail connect successive frozen problems. We give explicit common-contraction conditions for the scheduled dynamics using a Schur skeleton and modulation-derivative budgets; the measurement re-encoding residual remains an independent condition on the chosen lift and system. Under these contraction, residual, and reference conditions, a direct analysis of the implemented frozen optimal values yields practical exponential tracking. The bound separates freezing, re-encoding, reference defect, and certified feasible QP suboptimality, without requiring a second online nonlinear rollout or terminal optimization.
\end{abstract}
\begin{IEEEkeywords}
Koopman modeling, temporal Transformer, history-conditioned prediction, low-rank scheduling, convex model predictive control, practical exponential tracking.
\end{IEEEkeywords}

\section{Introduction}
\label{sec:introduction}
Model predictive control (MPC) selects control inputs by comparing their predicted effects over a finite horizon. Its prediction model must therefore capture the response of a nonlinear system to candidate inputs while remaining tractable for repeated online optimization. When the model is identified from input--output data, these requirements concern both the information used to initialize predictions and the structure used to propagate them \cite{ljung1999system,nelles2020nonlinear,forgione2021dynonet}. An informative representation is needed to distinguish different dynamical conditions, whereas a control-oriented evolution model is needed to evaluate their response to future inputs.

Koopman-based prediction offers a way to impose such an evolution structure: nonlinear observables provide lifted coordinates in which a matrix model approximates the dynamics. The underlying Koopman operator is linear but generally infinite-dimensional \cite{koopman1931hamiltonian,mezic2005spectral,brunton2022modern}; a computational predictor must instead use a finite set of observables. Extended dynamic mode decomposition (EDMD) estimates the corresponding finite-dimensional operator approximation from data \cite{williams2015data}. For control, the additional task is to represent the effect of externally chosen inputs. Dynamic mode decomposition with control separates state and input contributions, while controlled EDMD constructions yield lifted predictors with an explicit input channel \cite{proctor2016dmdc,korda2018linear}. With a linear input-affine predictor, quadratic costs and linear constraints can then lead to a convex MPC problem. This control advantage motivates the structure, but its usefulness still depends on finding suitable observables. Neural lifting methods address that choice by learning the coordinates together with their evolution rather than fixing the entire dictionary beforehand \cite{takeishi2017koopman,lusch2018deep,otto2019linearly,han2020deep}.

Learning the observables, however, cannot recover information that is absent from their inputs. If different physical states produce the same current measurement but different subsequent responses, a lift of that measurement alone cannot distinguish them. Past measurements and applied inputs can supply additional information about the preceding evolution. This motivates history-based representations: delay-coordinate approaches augment observables with past measurements \cite{kamb2020timedelay,das2019delay}, and recurrent or temporal autoencoders learn features from sequences for prediction and control \cite{han2024robust,shenoy2025temporal}. Transformer-based lifting extends sequence feature extraction to Koopman-inspired forecasting, as in DeepKoopFormer \cite{forootani2026deepkoopformer}. These approaches motivate using history to initialize the predictor; they do not require the physical state to be reconstructed explicitly.

History can also condition the evolution model itself. While a fixed transition matrix offers a simple prediction structure, allowing its coefficients to depend on recent data provides additional modeling flexibility. Koopa employs context-dependent operators for nonstationary forecasting \cite{liu2023koopa}, while MamKO generates matrices from historical data for modeling and predictive control \cite{li2025mamko}. These studies motivate distinguishing history-dependent initialization from history-conditioned state propagation. In receding-horizon prediction, however, the scheduling history must evolve with the predicted trajectory. The resulting coefficients depend on predicted measurements and candidate inputs, so conditional linearity at a fixed context does not preserve convexity of the full MPC problem.

To retain a quadratic program (QP), the coefficient sequence can be computed along a nominal rollout and held fixed while optimizing the candidate inputs. Quasi-linear parameter-varying MPC similarly exploits fixed scheduling sequences to form tractable subproblems, including iterative QP schemes \cite{cisneros2016qlmpc,cisneros2020qlpv}. Here we consider the specific policy of one nominal rollout and one frozen QP per sample, rather than iterating to update the schedules. The resulting approximation must be analyzed as part of the control law: at the next sample, new measurements change the lifted initialization and a new nominal rollout changes the prediction matrices. A decrease derived for one unchanged frozen model therefore does not establish a decrease across the problems actually solved. Likewise, results for real-time iteration concern a particular optimization update and cannot be transferred solely because both methods use a quadratic subproblem \cite{diehl2005realtime,gros2020rti}. The question is how the history-dependent model, the returned input sequence, and the next nominal rollout can be coupled to support a cross-sample tracking argument.

This paper develops a history-conditioned Transformer--Koopman predictor and a single-freeze MPC scheme. A temporal Transformer initializes the lifted state from measured history, while recursively updated measurement and input histories schedule low-rank transition corrections and measurement-dependent input modulation. At each sample, one nominal rollout determines the frozen linear time-varying prediction model used in a convex QP. The controller combines augmented-state tracking, reference-relative input increments, and a shifted input sequence with a reference tail. Under explicit common-contraction, re-encoding residual, and reference conditions, we establish practical exponential tracking by directly comparing the optimal values of consecutive frozen problems.

The main contributions are as follows:
\begin{enumerate}
\item \emph{A recursive history-conditioned predictor:} we combine a measurement-preserving Transformer lift with low-rank scheduling in a nonredundant augmented state, specifying how history initializes the features and conditions their subsequent evolution.
\item \emph{A single-freeze control law:} we couple one nominal rollout and one convex QP through a complete augmented tracking cost, reference-relative input increments, and a shift of the returned sequence with a reference tail.
\item \emph{Conditional closed-loop tracking analysis:} we construct sufficient common-contraction budgets and establish practical exponential tracking for the specified controller by comparing its successive frozen optimal values. The bound accounts separately for freezing, re-encoding, reference defects, and certified feasible QP suboptimality.
\end{enumerate}

Section~\ref{sec:prelim} formulates the prediction and control tasks. Section~\ref{sec:model} develops the model, learning objectives, and contractive design. Section~\ref{sec:control} presents the controller and its closed-loop analysis. Section~\ref{sec:experiments} reserves the numerical-results component, and Section~\ref{sec:conclusion} concludes. Fig.~\ref{fig:overview} summarizes the method.

\begin{figure*}[!t]
\centering
\setlength{\unitlength}{0.70mm}
\begin{picture}(220,127)
\linethickness{0.5pt}
\put(0,74){\framebox(45,25){\shortstack{Measured histories\\$\mX_k,\ v_k$}}}
\put(60,92){\framebox(51,20){\shortstack{Temporal Transformer\\$\vz_k^E=G_\theta(\mX_k)$}}}
\put(60,55){\framebox(51,20){\shortstack{Reconstructed history\\$\vc(q)=C_\phi(X(q),v)$}}}
\put(128,88){\framebox(45,24){\shortstack{Augmented state\\$q=\col(\vz,a,v)$}}}
\put(128,51){\framebox(45,24){\shortstack{Low-rank schedules\\$\mA(\vc),\ \mB(C_x\vz)$}}}
\put(185,69){\framebox(35,28){\shortstack{Recursive\\predictor\\$q^+=F(q,u)$}}}
\put(45,90){\rule{7\unitlength}{0.5pt}}\put(52,90){\rule{0.5pt}{12\unitlength}}\put(52,102){\makebox(8,0){$\rightarrow$}}
\put(22,99){\rule{0.5pt}{22\unitlength}}\put(22,121){\rule{128\unitlength}{0.5pt}}\put(150,112){\makebox(0,9){$\downarrow$}}
\put(86,124){\makebox(0,0){\scalebox{0.75}{\normalsize actual measurement and input caches}}}
\put(116,94){\rule{12\unitlength}{0.5pt}}\put(116,70){\rule{0.5pt}{24\unitlength}}\put(111,70){\makebox(5,0){$\leftarrow$}}
\put(111,102){\makebox(17,0){$\longrightarrow$}}\put(111,60){\makebox(17,0){$\longrightarrow$}}
\put(173,100){\rule{6\unitlength}{0.5pt}}\put(179,88){\rule{0.5pt}{12\unitlength}}\put(179,88){\makebox(6,0){$\rightarrow$}}
\put(173,63){\rule{6\unitlength}{0.5pt}}\put(179,63){\rule{0.5pt}{14\unitlength}}\put(179,77){\makebox(6,0){$\rightarrow$}}
\put(150,75){\makebox(0,13){$\downarrow$}}\put(165,81){\makebox(0,0){\scalebox{0.75}{\normalsize $C_x\vz$}}}
\put(202,44){\rule{0.5pt}{25\unitlength}}\put(23,44){\rule{179\unitlength}{0.5pt}}\put(23,36){\makebox(0,8){$\downarrow$}}
\put(111,47){\makebox(0,0){\scalebox{0.75}{\normalsize initialize at $q_k^E$; propagate once along $\bar U_k$}}}
\put(0,14){\framebox(46,22){\shortstack{One full nominal\\rollout of $F$}}}
\put(58,14){\framebox(46,22){\shortstack{Freeze once\\$\bar\mA_{i|k},\bar\mB_{i|k}$}}}
\put(116,14){\framebox(46,22){\shortstack{One convex QP\\return $\widehat U_k$}}}
\put(174,14){\framebox(46,22){\shortstack{Apply first input\\$\vu_k=\widehat\vu_{0|k}$}}}
\put(46,25){\makebox(12,0){$\longrightarrow$}}\put(104,25){\makebox(12,0){$\longrightarrow$}}\put(162,25){\makebox(12,0){$\longrightarrow$}}
\put(197,1){\rule{0.5pt}{13\unitlength}}\put(23,1){\rule{174\unitlength}{0.5pt}}\put(23,1){\makebox(0,13){$\uparrow$}}
\put(110,7){\makebox(0,0){\scalebox{0.75}{\normalsize next sample: re-encode measurements; shift $\widehat U_k$ and append $\vu_{k+N}^r$}}}
\end{picture}
\caption{Two roles of history and the single-freeze control cycle. The Transformer initializes the lifted features; recursively updated histories schedule subsequent transitions. The return arrow denotes the next sampling instant, not a second rollout within the current sample.}
\label{fig:overview}
\end{figure*}

\section{Preliminaries and Problem Formulation}
\label{sec:prelim}
\subsection{Koopman-Based Prediction}
\label{subsec:koopman_prelim}
Consider the discrete-time nonlinear system
\begin{equation}
\vs_{k+1}=f(\vs_k,\vu_k),\qquad \vx_k=h(\vs_k),
\label{eq:nonlinear_system}
\end{equation}
where $\vs_k\in\mathcal S\subseteq\R^{n_s}$ is the physical state and $\vu_k\in\R^m$ is the control input. The available measurement $\vx_k\in\R^n$ is obtained through the measurement map $h:\mathcal S\to\R^n$. The state transition map $f:\mathcal S\times\R^m\to\mathcal S$ governs the evolution of $\vs_k$ under $\vu_k$. We learn the predictor from measured input--output trajectories and use finite observation histories to capture temporal information not resolved by individual measurements.

For a fixed input $\vu$, the controlled Koopman operator acts on an observable $\varphi:\mathcal S\to\R$ as \cite{koopman1931hamiltonian,mezic2005spectral,brunton2022modern}
\begin{equation}
(\mathcal K_{\vu}\varphi)(\vs)=\varphi(f(\vs,\vu)).
\label{eq:koopman_operator}
\end{equation}
The operator is linear with respect to observables:
\begin{equation}
\mathcal K_{\vu}(\alpha\varphi+\beta\psi)=\alpha\mathcal K_{\vu}\varphi+\beta\mathcal K_{\vu}\psi.
\label{eq:koopman_linearity}
\end{equation}

Finite-dimensional Koopman models approximate this evolution using a selected set of observables. Since only $\vx=h(\vs)$ is measured, we construct the observables from $\vx$; each component $g_j(\vx)$ corresponds to the physical-state observable $g_j\circ h$. An instantaneous lifting map takes the form
\begin{equation}
\vz_k=g(\vx_k)=[g_1(\vx_k),\ldots,g_d(\vx_k)]^\top\in\R^d,
\label{eq:instantaneous_lifting}
\end{equation}
with the input-affine lifted approximation
\begin{equation}
\vz_{k+1}\approx\mA\vz_k+\mB\vu_k.
\label{eq:linear_lifted}
\end{equation}
Here $\mA\in\R^{d\times d}$ approximates the evolution of the selected observables, and $\mB\in\R^{d\times m}$ represents the control-input contribution. Together they define a finite-dimensional approximation of the controlled Koopman evolution in the chosen lifted coordinates.

To incorporate temporal information, we define the state-history window
\begin{equation}
\mX_k=\col(\vx_{k-P_h+1},\ldots,\vx_k)\in\R^{P_h n},\quad P_h\ge2,
\label{eq:measurement_window}
\end{equation}
and extend the lifting map to
\begin{equation}
\vz_k=G_\theta(\mX_k)=\col(\vx_k,\psi_\theta(\mX_k))\in\R^d,
\label{eq:state_preserving_lifting}
\end{equation}
where $d=n+n_\psi$. The lift retains the current measurement $\vx_k$ as its leading block, so the readout is fixed:
\begin{equation}
\vx_k=C_x\vz_k,\qquad C_x=[I_n\ 0].
\label{eq:state_readout}
\end{equation}
This representation incorporates recent measurements into the lifted state. The transition model may additionally depend on the recent trajectory context, as developed in Section~\ref{sec:model}.

\subsection{Problem Formulation}
\label{sec:problem}
The available training data comprise recorded measurement--input trajectories
\begin{equation}
\begin{aligned}
\mathcal D &= \{\tau^{(j)}\}_{j=1}^{N_{\rm tr}},\\
\tau^{(j)} &= \{(\vx_t^{(j)},\vu_t^{(j)})\}_{t=0}^{T_j},
\end{aligned}
\label{eq:training_data}
\end{equation}
where $N_{\rm tr}$ is the number of recorded trajectories and $T_j$ is the last time index of the $j$th trajectory.
For each admissible index $k$, the measurement history $\mX_k$ is defined in~\eqref{eq:measurement_window}, and the past-input history is
\begin{equation}
v_k = \col(\vu_{k-P_h+1},\ldots,\vu_{k-1}) \in \R^{(P_h-1)m}.
\label{eq:past_input_history}
\end{equation}
The current input $\vu_k$ is applied after the observed history window.

Given the measurement history $\mX_k$, the past-input history $v_k$, and a future control sequence $\vu_{k:k+H_p-1}$, the prediction task is to estimate $\vx_{k+1},\ldots,\vx_{k+H_p}$ over a horizon of $H_p$ steps. The predictor should exploit the observed history to initialize the lifted representation and support recursive multi-step prediction. The learned predictor is then used for receding-horizon reference tracking subject to input constraints $\vu_k\in\mathcal U$, where $\mathcal U=\{u:H_uu\leq h_u\}$ is a nonempty compact convex polytope.

The objective of this work is to develop a history-conditioned lifted predictor and a computationally tractable receding-horizon controller, and to establish a practical exponential tracking guarantee for the implemented closed loop under suitable conditions.

\section{History-Conditioned Scheduled Koopman Modeling}
\label{sec:model}
\subsection{History-Conditioned Koopman Predictor}
\label{sec:temporal_lifting}
Given the history window $\mX_k$ in~\eqref{eq:measurement_window}, we arrange its successive measurement blocks as the rows of a $P_h\times n$ array and partition it into $N_p=P_h/p$ contiguous non-overlapping temporal patches of length $p$, where $p$ is assumed to divide $P_h$. The $j$th patch is denoted by
\begin{equation}
\mX_k^{(j)}\in\R^{p\times n},\qquad j=1,\ldots,N_p.
\label{eq:temporal_patch}
\end{equation}
Each patch is vectorized as
\begin{equation}
t_k^{(j)}=\operatorname{vec}(\mX_k^{(j)})\in\R^{pn},
\label{eq:patch_vec}
\end{equation}
which is then mapped into the Transformer embedding space through a shared affine layer,
\begin{equation}
\vphi_k^{(j)}=\mW_p t_k^{(j)}+\vb_p\in\R^{d_{\rm model}},
\label{eq:patch_embedding}
\end{equation}
where $\mW_p\in\R^{d_{\rm model}\times pn}$ and $\vb_p\in\R^{d_{\rm model}}$. This patch-based tokenization reduces the sequence length from $P_h$ to $N_p=P_h/p$ while preserving local temporal structure within each patch.

The token sequence is augmented with sinusoidal positional encoding \cite{vaswani2017attention} to incorporate temporal order. With $\vp^{(j)}\in\R^{d_{\rm model}}$ the positional encoding of the $j$th patch, the augmented tokens are
\begin{equation}
\tilde{\vphi}_k^{(j)}=\vphi_k^{(j)}+\vp^{(j)},
\label{eq:pos_enc_applied}
\end{equation}
and the resulting token sequence is encoded by an $L_T$-layer Transformer encoder as
\begin{equation}
\begin{split}
\mH_k&=\operatorname{TransformerEncoder}
  \bigl(\{\tilde{\vphi}_k^{(j)}\}_{j=1}^{N_p}\bigr)\\
&\in\R^{N_p\times d_{\rm model}},
\end{split}
\label{eq:transformer_enc}
\end{equation}
where $L_T$ is the number of encoder layers and $\mH_k^{(j)}\in\R^{d_{\rm model}}$ denotes the contextualized representation of the $j$th patch.

We aggregate the patch features by mean pooling to obtain a fixed-dimensional history descriptor:
\begin{equation}
\bar\mH_k=\frac1{N_p}\sum_{j=1}^{N_p}\mH_k^{(j)}\in\R^{d_{\rm model}}.
\label{eq:pooling}
\end{equation}
After that, the pooled representation is projected to the history-feature space,
\begin{equation}
\psi_\theta(\mX_k)=\mW_h\bar\mH_k+\vb_h\in\R^{d-n},
\label{eq:history_feature}
\end{equation}
where $\mW_h\in\R^{(d-n)\times d_{\rm model}}$ and $\vb_h\in\R^{d-n}$.

The lifted state in~\eqref{eq:state_preserving_lifting} is thus
\begin{equation}
\vz_k=\col(\vx_k,\psi_\theta(\mX_k)),
\label{eq:lifted_state_form}
\end{equation}
where the first $n$ components are the current measurement $\vx_k=h(\vs_k)$ and the remaining $d-n$ components are the temporal features. This structure allows exact recovery of the measured coordinates through the fixed projection in~\eqref{eq:state_readout},
\begin{equation}
\vx_k=C_x\vz_k,\qquad C_x=[I_n\ 0],
\label{eq:decoder_projection}
\end{equation}
where $C_x$ extracts the measurement coordinates from the lifted state.

Although the temporal lifting provides different latent representations for different histories, a fixed transition $\mA$ still requires all such history-dependent latent states to evolve under one global linear operator. To relax this additional restriction without discarding the conditional-linear structure, we decompose the transition into a shared nominal component and a structured history-dependent correction. The input matrix is likewise allowed a measurement-dependent low-rank correction.

\newtheorem{definition}{Definition}
\begin{definition}[History-Conditioned Low-Rank Koopman Family]
\label{def:hcltk}
Let the history window be $\mX_k$ in~\eqref{eq:measurement_window}, and let $v_k$ denote the past-input history in~\eqref{eq:past_input_history}. Define the lifted state and context vector as
\begin{equation}
\begin{aligned}
\vz_k^E&=G_\theta(\mX_k)=\col(\vx_k,\psi_\theta(\mX_k)),\\
\vc_k&=C_\phi(\mX_k,v_k)\in\R^{d_c},
\end{aligned}
\label{eq:context_vector}
\end{equation}
where $G_\theta$ encodes the measurement history and $C_\phi:\R^{P_hn+(P_h-1)m}\to\R^{d_c}$ maps the measurement and past-input histories to the context vector, excluding the current input. The superscript $E$ denotes the encoder-generated lifted state, while $P$ denotes the model-propagated lifted state. The one-step prediction is
\begin{equation}
\vz_{k+1}^P=\mA(\vc_k)\vz_k^E+\mB(C_x\vz_k^E)\vu_k,
\label{eq:scheduled_predictor}
\end{equation}
with the low-rank schedules
\begin{align}
\mA(\vc)&=\mA_0+\mU_A\operatorname{diag}(\gamma_A(\vc))\mV_A^T,
\label{eq:transition_schedule}\\
\mB(\vx)&=\mB_0+\mU_B\operatorname{diag}(\gamma_B(\vx))\mV_B^T.
\label{eq:low_rank_schedules}
\end{align}
Here $\mA_0\in\R^{d\times d}$ and $\mB_0\in\R^{d\times m}$ are nominal matrices; $\mU_A,\mV_A\in\R^{d\times r_A}$, $\mU_B\in\R^{d\times r_B}$, and $\mV_B\in\R^{m\times r_B}$ are low-rank factors. The correction ranks are at most $r_A$ and $r_B$, respectively. The factors define correction directions, while the schedules determine their signed strengths. For $H\in\{A,B\}$, with argument $y=\vc$ or $y=\vx$, respectively, use
\begin{equation}
\gamma_{H,j}(y)=g_H[\operatorname{softmax}(W_H^\pi y+b_H^\pi)]_j
\tanh([W_Hy+b_H]_j).
\label{eq:A_modulation_coefficients}
\end{equation}
\end{definition}
Unit-bounded factor columns give $\norm{\mA-\mA_0}_2\le g_A$ and $\norm{\mB-\mB_0}_2\le g_B$. This bounds the correction magnitude, not the stability of the scheduled dynamics. Fixed scheduling signals give conditional linear propagation in $(\vz,\vu)$; recursive scheduling is generally nonlinear.

In recursive prediction, the encoder initializes the lifted state once, while subsequent steps propagate $\vz^P$ with the context updated from predicted histories. To express this recursion without duplicating the current measurement, introduce the measurement cache and augmented state
\begin{equation}
\begin{split}
a_k&=\col(\vx_{k-P_h+1},\ldots,\vx_{k-1}),\\
q_k&=\col(\vz_k,a_k,v_k)\in\R^{s_q},\quad
s_q=d+(P_h-1)(n+m).
\end{split}
\label{eq:augmented_state}
\end{equation}
Let $E_{\rm prev}\in\R^{m\times s_q}$ select the last input slot of $v$. Let $S_x,S_u$ shift the respective caches, and $E_x,E_u$ append one measurement or input. The complete predictor is
\begin{equation}
\begin{split}
\vc(q)&=C_\phi(\col(a,C_x\vz),v),\\
F(q,\vu)&=\col\bigl(\mA(\vc(q))\vz+\mB(C_x\vz)\vu,\\
&\hspace{19mm}S_xa+E_xC_x\vz,\ S_uv+E_u\vu\bigr).
\end{split}
\label{eq:complete_control_model}
\end{equation}
Thus $a^+$ appends the \emph{current} $C_x\vz$ and $v^+$ appends the applied input. Combining $a^+$ with $C_x\vz^+$ recovers the next full measurement window, equivalent to shifting the window by the next predicted measurement while avoiding duplication of the current measurement in the propagated state. Writing $X(q)=\col(a,C_x\vz)$, the upper branches of Fig.~\ref{fig:overview} distinguish initial encoding from recursive history conditioning.

At a control sample, $\vz_k^E=G_\theta(\mX_k)$ is recomputed from measurements and combined with the actual caches. Subsequent predictions propagate $q$ using~\eqref{eq:complete_control_model} without re-encoding. Their discrepancy from the next measured initialization is
\begin{equation}
q_{k+1}^E=F(q_k^E,\vu_k)+w_k.
\label{eq:residual_definition}
\end{equation}
With exact cache bookkeeping and input execution, $w_k=\col(r_k,0,0)$. The lifted residual $r_k$ includes both measurement prediction error and the defect between recursively propagated features and the encoder image. Appendix~\ref{app:construction} makes this decomposition explicit; a small training loss alone does not establish a uniform bound on $w_k$.

\subsection{Model Learning and Contractive Design}
\label{sec:training}
For a positive-definite matrix $\Pi\succ0$, we write $\norm{e}_\Pi^2=e^\top\Pi e$ for the weighted squared norm of a vector $e$, and $\norm{L}_\Pi=\norm{\Pi^{1/2}L\Pi^{-1/2}}_2$ for the induced norm of a square linear map $L$.

We formulate learning within a constrained family, with prediction losses fitting the data and structural bounds controlling recursive sensitivity. These roles are distinct: the losses couple the encoder to the propagated features, whereas the constraints below select the scheduled dynamics used by the controller. We first initialize the nominal predictor using one-step targets and then fine-tune the scheduled model through recursive rollouts. For measured windows define $\vz_t^E=G_\theta(\mX_t)$. At each update the target branch is recomputed and held fixed, denoted by $\operatorname{sg}(\cdot)$; the predictive branch remains trainable. The initialization stage fits the lift and nominal dynamics using $\widetilde\vz_{k+1}^0=\mA_0G_\theta(\mX_k)+\mB_0\vu_k$ and
\begin{equation}
\begin{split}
\mathcal L_{\rm init}=\sum_{k\in\mathcal I_1}\bigl[&
\norm{C_x\widetilde\vz_{k+1}^0-\vx_{k+1}}_2^2\\
&+\alpha_0\norm{\widetilde\vz_{k+1}^0-\operatorname{sg}(\vz_{k+1}^E)}_2^2\bigr],
\end{split}
\label{eq:training_one_step_objects}
\end{equation}
where $\alpha_0>0$ and $\mathcal I_1$ contains valid one-step pairs. The measured-output term directly supervises the output rows; the encoded-state term supplies an additional consistency objective rather than an identification theorem for all latent dynamics.

The scheduled-rollout stage activates both schedules. From each observed initial history, set $q_{k|k}^P=q_k^E$ and generate $q_{k+i+1|k}^P=F(q_{k+i|k}^P,\vu_{k+i})$ for a training horizon $H_R$. Writing $\vz^P=E_z^Tq^P$, where $E_z$ injects the lifted coordinates into $q$, minimize
\begin{equation}
\begin{split}
\mathcal L_R=\sum_{k\in\mathcal I_R}\sum_{i=1}^{H_R}\omega_i\bigl[&
\norm{C_x\vz_{k+i|k}^P-\vx_{k+i}}_2^2\\
&+\alpha_R\norm{\vz_{k+i|k}^P-\operatorname{sg}(\vz_{k+i}^E)}_2^2\bigr],
\end{split}
\label{eq:joint_training_loss}
\end{equation}
with $\omega_i=1/H_R$ and $\alpha_R>0$. Samples and target windows come from the same recorded trajectory. The encoder initializes each predictive rollout once; the context then follows predicted measurement histories and prescribed inputs. The encoder, context, and allowable dynamics parameters are trained jointly, with $C_x$ fixed.

Repeated prediction makes the sensitivity of both the state update and its history schedule relevant to control. A small empirical rollout loss does not bound that sensitivity uniformly. We therefore restrict the dynamics through a stable skeleton and explicit modulation budgets, rather than treating the control assumptions as consequences of fitting. Choose a Schur $\mA_0$ and define
\begin{equation}
M_0=\begin{bmatrix}\mA_0&0&0\\E_xC_x&S_x&0\\0&0&S_u\end{bmatrix},\qquad
M_0^T\Pi M_0-\Pi=-I.
\label{eq:lyapunov_skeleton}
\end{equation}
The cache shifts are nilpotent, so $M_0$ is Schur and $\Pi\succ0$ exists. Schur stability of $\mA_0$ is a design restriction; $\Pi$ is the analysis norm for that fixed skeleton. One can train $\mA_0=TKT^{-1}$ with $T$ invertible and $\norm{K}_2<1$, or keep it fixed during fine-tuning. Any change of the skeleton requires a new $\Pi$ and new gain budgets before deployment.

On the ball $\mathcal Q=\{q:\norm{q}_\Pi\le M\}$, let $L_c$ bound the Euclidean derivative of $q\mapsto\vc(q)$. For each head define
\begin{equation}
\kappa_H=\norm{W_H}_2+\tfrac12\sqrt{r_H}\norm{W_H^\pi}_2,\qquad L_H=g_H\kappa_H.
\label{eq:Lipschitz_constants}
\end{equation}
Factor and head-norm constraints give explicit amplitude and derivative bounds, with sufficient MLP and attention estimates in Appendix~\ref{app:construction}. The resulting construction restricts the \emph{scheduled dynamics}; it does not establish a uniform encoder residual for every Transformer. That residual, which includes the encoding-graph defect, remains a separate assumption on the chosen lift and measured system. Training must preserve the structural bounds, or its final fixed parameters must be checked against them. Deployment is deterministic, with dropout disabled and normalization fixed.

\begin{proposition}[Constructible contractive model class]
\label{prop:contractive_class}
Fix a Schur skeleton \eqref{eq:lyapunov_skeleton}, a compact input polytope, finite head and context derivative bounds, and an independently available residual budget $\bar w$ on a trial ball $\mathcal Q$. Let $\lambda_0=\norm{M_0}_\Pi<1$ and let $b_0^{\rm drv}$ bound the nominal input drive. If the ball covers the initial set and
\begin{equation}
\Delta=(1-\lambda_0)M-b_0^{\rm drv}-\bar w>0,
\label{eq:radius_margin}
\end{equation}
then the explicit positive-gain choices in \eqref{eq:gain_choice} yield uniform same-input contraction of the complete model and every frozen model generated by a nominal point in $\mathcal Q$. Their rates $\lambda,\lambda_{\rm fr}<1$ and drive bound $b_u$ satisfy
\begin{equation}
\lambda M+b_u+\bar w<M,\qquad \lambda_{\rm fr}M+b_u<M.
\label{eq:invariant_radius}
\end{equation}
Both modulation gains can be nonzero. Nonconstant heads and effective factor directions give genuinely history-dependent members of the class.
\end{proposition}
\begin{proof}
Schur stability gives the common Lyapunov norm. The full derivative perturbation includes the direct correction, the derivative of the history schedule acting on $\vz$, and the derivative of $\mB(C_x\vz)\vu$. Their bounds give $\norm{D_qF-M_0}_\Pi\le\nu_{\rm bd}$; the frozen bound omits the schedule derivatives. Fixing the ball first makes these perturbation bounds linear in the common gain scale. The strict choice \eqref{eq:gain_choice} leaves positive contraction and radius margins. Applying the radius inequalities to a currently admissible point proves next-step membership, rather than assuming future trajectories remain in the ball. The detailed bounds and a nonconstant, encoding-consistent periodic example are given in Appendix~\ref{app:construction}.
\end{proof}

The proposition constructs a nonempty constrained model class; it does not ensure that an arbitrary training result belongs to it. Compatibility of these structural bounds with a small uniform residual remains an independent condition on the model and measured system. On the same domain, contraction also gives the finite-horizon open-loop estimate $\norm{q_{k+i}^E-q_{k+i|k}^P}_\Pi\le\bar w\sum_{j=0}^{i-1}\lambda^j$ for a common initial history and identical prescribed inputs.

\section{Frozen-LTV Model Predictive Control}
\label{sec:control}
\subsection{Frozen-LTV MPC Formulation}
\label{subsec:frozen_mpc}
Fix a receding horizon $N\ge1$. A tracking reference is a fixed absolute-time sequence $\{(q_t^r,\vu_t^r)\}_{t\ge0}$ in the augmented state coordinates defined in Section~\ref{sec:model}, with $q_t^r\in\mathcal Q$ and $\vu_t^r\in\mathcal U$ for every required time $t$.

The scheduled predictor is nonlinear in a candidate control sequence because future schedules depend on predicted histories. To preserve a convex online problem, initialize $q_k^E$ from the current measurement history and select a feasible nominal sequence $\bar U_k\in\mathcal U^N$. Perform one complete rollout,
\begin{equation}
\bar q_{0|k}=q_k^E,\qquad
\bar q_{i+1|k}=F(\bar q_{i|k},\bar\vu_{i|k}),\quad 0\le i<N,
\label{eq:nominal_rollout}
\end{equation}
and record $\bar\mA_{i|k}=\mA(\vc(\bar q_{i|k}))$ and $\bar\mB_{i|k}=\mB(C_xE_z^T\bar q_{i|k})$. Hold these coefficients fixed throughout the optimization. The resulting candidate trajectory obeys
\begin{equation}
\begin{split}
p_{0|k}&=q_k^E,\qquad p_{i+1|k}=M_{i|k}p_{i|k}+D_{i|k}\vu_{i|k},\\
M_{i|k}&=\begin{bmatrix}\bar\mA_{i|k}&0&0\\E_xC_x&S_x&0\\0&0&S_u\end{bmatrix},\quad
D_{i|k}=\col(\bar\mB_{i|k},0,E_u).
\end{split}
\label{eq:frozen_model}
\end{equation}
The coefficients come from the nominal trajectory, not the decision trajectory. Therefore $p_{i|k}$ is affine in the stacked decision $U$, even though the complete scheduled trajectory generally is not. The lower branch of Fig.~\ref{fig:overview} shows this separation between generating the coefficients and optimizing the candidate inputs.

For fixed absolute-time references, use
\begin{equation}
\begin{split}
\ell_t(q,u)={}&\norm{q-q_t^r}_\Pi^2+\norm{u-\vu_t^r}_R^2\\
&+\norm{(u-E_{\rm prev}q)-(\vu_t^r-\vu_{t-1}^r)}_S^2,
\end{split}
\label{eq:stage_cost}
\end{equation}
where $R,S\succeq0$. Penalizing the complete augmented error provides the positive-definite state term needed by the value analysis. The relative increment penalizes departures from the reference's own input change: applying the reference tail leaves only the previous-input cache error, including for a time-varying reference.

With terminal weight $\beta\Pi$, solve the single problem
\begin{equation}
\begin{split}
W_k&=\min_{U\in\mathcal U^N}J_k^{\rm fr}(U;\bar U_k),\\
J_k^{\rm fr}&=\sum_{i=0}^{N-1}\ell_{k+i}(p_{i|k},\vu_{i|k})
+\beta\norm{p_{N|k}-q_{k+N}^r}_\Pi^2.
\end{split}
\label{eq:quadratic_template}
\end{equation}
Choose $\beta$ by \eqref{eq:terminal_weight} below, using the common contraction rate. Only the input polytope is a hard constraint; there is no state constraint, terminal set, or online terminal optimization.

The frozen objective is a sum of positive-semidefinite quadratic terms in affine functions of $U$, and the constraints are $(I_N\otimes H_u)U\le\mathbf1_N\otimes h_u$. A minimizer exists by continuity and compactness. Its Hessian is in fact positive definite: the first nonzero input variation enters the next input-cache slot, which is penalized by a complete-state stage weight or the terminal weight. This is specific to the augmented cost and holds even when $R,S$ are singular. The condensed expression and proof are in Appendix~\ref{app:estimates}.

Let $\widehat U_k$ be the returned feasible sequence. Apply and shift that same sequence:
\begin{align}
\vu_k&=\widehat\vu_{0|k},\label{eq:apply_first}\\
\bar U_{k+1}&=(\widehat\vu_{1|k},\ldots,
\widehat\vu_{N-1|k},\vu_{k+N}^r).
\label{eq:nominal_shift}
\end{align}
For an exact solve $\widehat U_k=U_k^*$; a feasible numerical solve is covered by an objective-gap condition below. When $N=1$, \eqref{eq:nominal_shift} contains only the reference tail. This nominal update is part of the control law, not an arbitrary warm-start choice. The complete online cycle is summarized in Algorithm~\ref{alg:mpc}.

\begin{algorithm}[!t]
\caption{Single-Freeze Transformer--Koopman MPC}
\label{alg:mpc}
\begin{algorithmic}[1]
\STATE Fix the deployed model, $\Pi$, $\mathcal U$, $N$, $\beta$, and absolute-time reference; initialize any $\bar U_0\in\mathcal U^N$.
\FOR{$k=0,1,\ldots$}
\STATE Measure $\mX_k$, re-encode $\vz_k^E$, and assemble $q_k^E$ using actual histories.
\STATE Perform one full nominal rollout \eqref{eq:nominal_rollout} and freeze its coefficients in \eqref{eq:frozen_model}.
\STATE Solve one QP \eqref{eq:quadratic_template}; return feasible $\widehat U_k$ with the required objective-gap bound.
\STATE Apply \eqref{eq:apply_first} and update the next nominal by \eqref{eq:nominal_shift}.
\ENDFOR
\end{algorithmic}
\end{algorithm}

\subsection{Closed-Loop Stability Analysis}
\label{subsec:closed_loop}
The main difficulty is that the next sample changes both the measured initialization and the coefficients of the QP. A decrease proved for an unchanged frozen model would therefore miss the implemented controller. We compare the optimal values of the \emph{successive re-frozen problems}, using the nominal shift in \eqref{eq:nominal_shift} to construct the next feasible candidate. The value $W_k$ depends on $q_k^E$, $\bar U_k$, and absolute reference time, rather than on the state alone.

\begin{assumption}[Common model and residual domain]
\label{ass:domain}
Fix the deployed model, $\Pi\succ0$, $M>0$, and $\mathcal Q=\{q:\norm{q}_\Pi\le M\}$, with $q_0^E\in\mathcal Q$. Let
\begin{equation}
\mathcal S_{\mathcal Q}=\{(\mA(\vc(\bar q)),\mB(C_xE_z^T\bar q)):
\bar q\in\mathcal Q\}.
\label{eq:frozen_family}
\end{equation}
For every $q,q'\in\mathcal Q$, $u\in\mathcal U$, and each coefficient pair in $\mathcal S_{\mathcal Q}$, the complete and frozen maps satisfy
\begin{align}
\norm{F(q,u)-F(q',u)}_\Pi&\le\lambda\norm{q-q'}_\Pi,\nonumber\\
\norm{F^{\rm fr}(q,u)-F^{\rm fr}(q',u)}_\Pi
&\le\lambda_{\rm fr}\norm{q-q'}_\Pi,\label{eq:common_contraction}\\
\norm{F(0,u)}_\Pi,\ \norm{F^{\rm fr}(0,u)}_\Pi&\le b_u,\nonumber\\
\norm{F(q,u)-F^{\rm fr}(q,u)}_\Pi&\le\bar\eta,\nonumber
\end{align}
where $0\le\lambda,\lambda_{\rm fr}<1$. For every admissible measured transition whose current encoded state lies in $\mathcal Q$, the residual in \eqref{eq:residual_definition} satisfies $\norm{w_k}_\Pi\le\bar w$. Moreover,
\begin{equation}
\lambda M+b_u+\bar w\le M,\qquad
\lambda_{\rm fr}M+b_u\le M.
\label{eq:domain_radius}
\end{equation}
The residual bound does not presuppose next-state membership in $\mathcal Q$. All constants are uniform in time, nominal sequences, and candidate inputs.
\end{assumption}

Proposition~\ref{prop:contractive_class} supplies sufficient structural choices for the common domain; the actual residual bound remains independent. The amplitude and norm bounds in Appendix~\ref{app:construction} also provide a uniform choice of $\bar\eta$. In particular, the assumption covers all nominal coefficient pairs generated in the ball, not only those near the reference.

\begin{assumption}[Reference consistency]
\label{ass:reference}
The fixed absolute-time reference satisfies, for every required time $t$,
\begin{equation}
\begin{split}
q_t^r\in\mathcal Q,\quad \vu_t^r\in\mathcal U,\quad
E_{\rm prev}q_t^r=\vu_{t-1}^r,\\
\norm{F(q_t^r,\vu_t^r)-q_{t+1}^r}_\Pi\le\bar d_{\rm ref}.
\end{split}
\label{eq:reference_defect}
\end{equation}
These sequences are not revised at a later sample for the same absolute time.
\end{assumption}

Let $H_S=S^{1/2}E_{\rm prev}\Pi^{-1/2}$, $c_S=1+\norm{H_S}_2^2$. Choose $\vartheta>0$ with $(1+\vartheta)\lambda^2<1$ and
\begin{equation}
\beta\ge\frac{c_S}{1-(1+\vartheta)\lambda^2}.
\label{eq:terminal_weight}
\end{equation}
This choice makes the reference input a suitable tail: its relative-increment cost contains only the previous-input error, while contraction controls the next terminal error. The shift budget $\varepsilon$ is defined in \eqref{eq:epsilon_closed_loop} of Appendix~\ref{app:estimates}; the uniform value-comparison constants $\mathsf A,\mathsf B$ are specified in \eqref{eq:value_constants} of Appendix~\ref{app:comparison}. These constants quantify the effects of re-encoding, freezing, and reference defect in the following result.

\begin{theorem}[Practical exponential tracking of single-freeze MPC]
\label{thm:tracking}
Under Assumptions~\ref{ass:domain}--\ref{ass:reference}, use \eqref{eq:terminal_weight} and execute Algorithm~\ref{alg:mpc} with exact input application. Suppose the returned sequence is primal feasible and satisfies
\begin{equation}
0\le J_k^{\rm fr}(\widehat U_k;\bar U_k)-W_k
\le\delta_k\le\bar\delta,
\label{eq:solver_gap}
\end{equation}
with $\delta_k=0$ for an exact solve. All actual and predicted augmented states remain in $\mathcal Q$, and input feasibility is recursive. With the constants defined above, $e_k=q_k^E-q_k^r$, and $\rho=1-1/\mathsf A\in[0,1)$,
\begin{equation}
\begin{split}
\norm{e_k}_\Pi^2&\le W_k\le\Xi_k,\\
\Xi_k&=\mathsf A\rho^k\norm{e_0}_\Pi^2+\mathsf B\\
&\quad+\sum_{t=0}^{k-1}\rho^{k-1-t}(\varepsilon+\delta_t).
\end{split}
\label{eq:squared_tracking_bound}
\end{equation}
In particular,
\begin{equation}
\norm{e_k}_\Pi\le\sqrt{\mathsf A}\rho^{k/2}\norm{e_0}_\Pi
+\sqrt{\mathsf B+\mathsf A(\varepsilon+\bar\delta)},
\label{eq:practical_tracking}
\end{equation}
and $\limsup_{k\to\infty}\norm{e_k}_\Pi\le
\sqrt{\mathsf B+\mathsf A(\varepsilon+\bar\delta)}$.
\end{theorem}
\begin{proof}
For a current point in $\mathcal Q$, the complete next-state norm is at most $\lambda M+b_u$, and the frozen bound is $\lambda_{\rm fr}M+b_u$. Equation~\eqref{eq:domain_radius} first keeps the nominal rollout in the domain, so all its coefficients are admissible; induction then keeps every complete and frozen candidate there. Adding $\bar w$ keeps the next actual point in the same ball. The input feasible set is always $\mathcal U^N$, and the shifted sequence is feasible because its tail lies in $\mathcal U$.

For any candidate first input, the full and frozen first updates coincide: both use coefficients evaluated at the common initial state $q_k^E$. By construction, the frozen prediction coincides with the full scheduled rollout along the nominal input sequence. Induction establishes this identity for that sequence, not full-horizon equality for arbitrary candidate inputs.

Suppress the old sample index and let $p_i$ be the old frozen trajectory of $\widehat U_k$. First-step exactness gives $q_{k+1}^E=p_1+w_k$. Let $U^+$ be \eqref{eq:nominal_shift} and $\widetilde q_j$ the full trajectory from $q_{k+1}^E$ along $U^+$. The initial distance from $p_1$ is at most $d_0=\bar w$; the intermediate bounds satisfy $d_{j+1}=\lambda d_j+\bar\eta$ for $0\le j<N-1$. With the explicit $d_j$ in \eqref{eq:shift_error}, this gives
\begin{equation}
\begin{split}
\norm{\widetilde q_j-p_{j+1}}_\Pi&\le d_j\quad(0\le j<N),\\
\norm{\widetilde q_N-F(p_N,\vu_{k+N}^r)}_\Pi
&\le\lambda d_{N-1}.
\end{split}
\label{eq:shift_comparison}
\end{equation}
The last comparison uses the same full map and reference input on both sides, so there is no additional frozen mismatch at the tail.

The tail input also gives a cost decrease. Its input-error term vanishes and its relative-increment cost is at most $(c_S-1)\norm{q-q_t^r}_\Pi^2$. Contraction and Young's inequality, together with \eqref{eq:terminal_weight}, imply
\begin{equation}
\begin{split}
\ell_t(q,\vu_t^r)+\beta\norm{F(q,\vu_t^r)-q_{t+1}^r}_\Pi^2
-\beta\norm{q-q_t^r}_\Pi^2\le\tau_{\rm ref},
\end{split}
\label{eq:tail_decrease}
\end{equation}
with $\tau_{\rm ref}$ defined in Appendix~\ref{app:estimates}. Because the next nominal is exactly $U^+$, the new frozen candidate trajectory equals the full rollout $\widetilde q$. At matched absolute times and inputs, the cost continuity bounds \eqref{eq:cost_continuity} convert the trajectory differences into the budget $\varepsilon$ of \eqref{eq:epsilon_closed_loop}. Applying the tail inequality \eqref{eq:tail_decrease} at the old endpoint $p_N$ then yields
\begin{align}
W_{k+1}&\le J_{k+1}^{\rm fr}(U^+;U^+)\nonumber\\
&\le J_k^{\rm fr}(\widehat U_k;\bar U_k)
-\ell_k(q_k^E,\widehat\vu_{0|k})+\varepsilon\nonumber\\
&\le W_k-\norm{e_k}_\Pi^2+\varepsilon+\delta_k.
\label{eq:value_descent}
\end{align}
This compares the actual re-frozen values, rather than optimal values of an ideal nonlinear MPC. The full scheduled trajectory $\widetilde q$ is used only in the analysis and is not computed in addition to the next sample's nominal rollout.

The initial stage gives $W_k\ge\norm{e_k}_\Pi^2$. For the upper bound choose the feasible reference input sequence $(\vu_k^r,\ldots,\vu_{k+N-1}^r)$. Contraction bounds its full-state reference error, and \eqref{eq:frozen_discrepancy} bounds the frozen-model discrepancy by $D_i$. Thus its state error is at most $\lambda^i\norm{e_k}_\Pi+b_i$, with $b_i$ given in Appendix~\ref{app:comparison}. Its $R$ terms vanish; after step zero, its previous-input slot equals the preceding reference input, so its relative-increment terms also vanish. Squaring and summing gives
\begin{equation}
\norm{e_k}_\Pi^2\le W_k\le\mathsf A\norm{e_k}_\Pi^2+\mathsf B.
\label{eq:value_comparison}
\end{equation}
Substitute this upper comparison into \eqref{eq:value_descent} to obtain $W_{k+1}\le\rho W_k+\mathsf B/\mathsf A+\varepsilon+\delta_k$. Expanding the scalar recursion and using $W_0\le\mathsf A\norm{e_0}_\Pi^2+\mathsf B$ gives $W_k\le\Xi_k$. Combining this with the lower comparison proves \eqref{eq:squared_tracking_bound}. The geometric sum is at most $\mathsf A$; subadditivity of the square root gives \eqref{eq:practical_tracking} and its limiting radius.
\end{proof}

The radius separates the model/reference comparison error $\mathsf B$, the cross-sample shift budget $\varepsilon$, and the feasible solver gap $\bar\delta$. All caches are included in the augmented norm; general cache perturbations need not vanish. Empty sums are zero and $\lambda^0=1$. In particular, for $N=1$ the intermediate shift recursion is empty and $\varepsilon=\tau_{\rm ref}+(L_\ell+\lambda L_f)\bar w$. For an exact model reference one can use the sharper $\beta\ge c_S/(1-\lambda^2)$ with $\tau_{\rm ref}=0$.

\begin{corollary}[Measured and input tracking]
\label{cor:projections}
Let $C_m=C_xE_z^T$ and use $\Xi_k$ from \eqref{eq:squared_tracking_bound}. Then
\begin{equation}
\norm{\vx_k-C_mq_k^r}_2\le\norm{C_m\Pi^{-1/2}}_2\norm{e_k}_\Pi.
\label{eq:measurement_tracking}
\end{equation}
If $R\succ0$, then
\begin{equation}
\norm{\widehat\vu_{0|k}-\vu_k^r}_2^2\le
\frac{\Xi_k+\delta_k}{\lambda_{\min}(R)}.
\label{eq:input_tracking}
\end{equation}
\end{corollary}
\begin{proof}
The first bound is a fixed projection because the lift preserves the actual measurement. If $R\succ0$, the squared input error times $\lambda_{\min}(R)$ is at most the delivered cost. The gap condition and the value bound in \eqref{eq:squared_tracking_bound} give $J_k^{\rm fr}(\widehat U_k;\bar U_k)\le W_k+\delta_k\le\Xi_k+\delta_k$, proving \eqref{eq:input_tracking} with the current objective gap retained.
\end{proof}

\begin{remark}
If $\bar w=\bar\eta=\bar d_{\rm ref}=\bar\delta=0$, the estimate gives exponential tracking to zero. For $N\ge2$, vanishing re-encoding and reference defects alone do not generally eliminate freezing-induced tracking error. The conclusion concerns the encoded augmented state and its measurement projection, not unobserved physical-state stability, general input-to-state stability, or state and obstacle safety. Common same-input contraction is a sufficient model-class restriction; its compatibility with a small physical prediction residual must be assessed independently, and neither a convex QP nor this theorem establishes a hardware execution deadline.
\end{remark}

\section{Numerical Results}
\label{sec:experiments}
Numerical results are not included in this manuscript version. This section retains their place in the complete paper; no empirical accuracy, comparison, hardware-validation, or timing claim is made. The modeling and closed-loop statements above are analytical and conditional on the stated assumptions.

\section{Conclusion}
\label{sec:conclusion}
This paper developed a history-conditioned scheduled Koopman predictor with a measurement-preserving temporal Transformer lift and bounded low-rank transition and input modulation. Nonredundant measurement and input caches define the recursively propagated augmented state. A Schur skeleton and explicit schedule-derivative budgets provide a constructible common-contraction model class without removing nonconstant history dependence.

For control, one nominal rollout and one frozen convex QP are coupled by a reference-relative augmented tracking objective and a shift of the returned sequence with a reference tail. A direct comparison of successive frozen optimal values establishes practical exponential tracking under uniform contraction, re-encoding residual, reference-defect, and feasible objective-gap conditions. The resulting radius separates these error contributions and projects directly to measured coordinates. The guarantee is for the specified single-freeze controller and model class; it is not an unconditional claim about an arbitrary learned predictor or hidden physical state.

\appendices
\section{Contractive Predictor Construction}
\label{app:construction}
The skeleton in \eqref{eq:lyapunov_skeleton} is block triangular, with eigenvalues those of $\mA_0$ and the nilpotent shift blocks. Thus it is Schur, and $\Pi=\sum_{j=0}^{\infty}(M_0^T)^jM_0^j$ is positive definite. Its induced norm satisfies $\lambda_0^2=1-1/\lambda_{\max}(\Pi)<1$.

For completeness, the softmax Jacobian $J(\pi)=\operatorname{diag}(\pi)-\pi\pi^T$ has Euclidean norm at most $1/2$. Combining this bound, the $1$-Lipschitz property of $\tanh$, and $\norm{b}_1\le\sqrt r\norm{b}_2$ gives \eqref{eq:Lipschitz_constants}: split the change of $\pi_j\tanh([Wy+b]_j)$ into the activation change and the softmax-weight change, then sum the unit-bounded rank-one factors. The first sum is at most $\norm{W}_2\norm{\Delta y}_2$, the second at most $\sqrt r\norm{W^\pi}_2\norm{\Delta y}_2/2$.

Let $L_c$ bound the context derivative on the selected ball. A two-layer $\tanh$ context has $L_c\le\norm{W_{c,2}}_2\norm{W_{c,1}}_2$. For a single attention head $Y(Z)=P(Z)V(Z)W_O$, $P=\operatorname{softmax}_{\rm row}(Q(Z)K(Z)^T/\sqrt{d_h})$, suppose the affine query, key, and value maps have Frobenius amplitude bounds $M_Q,M_K,M_V$ and Lipschitz constants $L_Q,L_K,L_V$ on a convex token domain. With $L_{\rm tok}$ tokens,
\begin{equation}
\begin{split}
L_{\rm score}&\le(L_QM_K+M_QL_K)/\sqrt{d_h},\\
L_Y&\le\norm{W_O}_2\left(\tfrac12M_VL_{\rm score}
+\sqrt{L_{\rm tok}}L_V\right).
\end{split}
\label{eq:attention_bound}
\end{equation}
This follows from the product rule, the rowwise softmax bound, and $\norm{P}_2\le\norm{P}_F\le\sqrt{L_{\rm tok}}$. Multihead concatenation combines the head bounds in Euclidean norm before the final projection. Residual blocks have constant at most one plus the branch constant; feedforward and embedding constants multiply under composition. Layer normalization with fixed positive $\epsilon_{\rm LN}$ has derivative norm at most $\norm{\gamma_{\rm LN}}_\infty/\sqrt{\epsilon_{\rm LN}}$, and mean pooling has norm $1/\sqrt{L_{\rm tok}}$. Fixed positions and biases enter amplitude bounds. These formulas permit constrained Transformer context maps, but do not imply that an unconstrained trained network has a useful small bound. When the temporal encoder uses GELU activations, its feature-map bound additionally incorporates a uniform activation-derivative bound.

Define the norm-conversion factors and domain bounds
\begin{equation}
\begin{split}
a_z&=\norm{\Pi^{1/2}E_z}_2,\quad
b_z=\norm{E_z^T\Pi^{-1/2}}_2,\\
b_q&=\norm{\Pi^{-1/2}}_2,\quad
b_x=\norm{C_xE_z^T\Pi^{-1/2}}_2,\\
M_u&=\sup_{u\in\mathcal U}\norm{u}_2,\qquad M_z=b_zM.
\end{split}
\label{eq:norm_conversion}
\end{equation}
For a variation $\Delta q$, with $\Delta\vz=E_z^T\Delta q$, the top block of $D_qF-M_0$ is the sum
\begin{equation}
\begin{split}
&(\mA(\vc)-\mA_0)\Delta\vz
+D\mA(\vc)[D\vc(q)\Delta q]\vz\\
&\hspace{14mm}+D\mB(C_x\vz)[C_x\Delta\vz]u.
\end{split}
\label{eq:full_derivative}
\end{equation}
The other blocks vanish. Therefore
\begin{equation}
\nu_{\rm bd}=a_z(g_A b_z+L_A L_cM_zb_q+L_BM_ub_x)
\label{eq:derivative_budget}
\end{equation}
bounds its $\Pi$-norm. The frozen perturbation is bounded by $a_zg_A b_z$, uniformly over all coefficients generated in the ball. Smoothness on the ball and its line segments turns these derivative bounds into same-input Lipschitz bounds.

To give a fully specified gain choice, let
\begin{equation*}
b_0^{\rm drv}=\sup_{u\in\mathcal U}
\norm{\Pi^{1/2}\col(\mB_0u,0,E_uu)}_2.
\end{equation*}
Choose a trial ball satisfying \eqref{eq:radius_margin}, and fix head weights and $L_c$ there. Take $g_A=t\widehat g_A$, $g_B=t\widehat g_B$ with $\widehat g_A,\widehat g_B>0$, and set
\begin{equation}
\begin{split}
k_F&=a_z(\widehat g_A b_z+\widehat g_A\kappa_A L_c b_zMb_q
+\widehat g_B\kappa_BM_ub_x),\\
k_{\rm fr}&=a_z\widehat g_A b_z,\quad k_b=a_z\widehat g_BM_u,
\quad k_* =\max(k_F,k_{\rm fr}).
\end{split}
\label{eq:gain_coefficients}
\end{equation}
Any choice
\begin{equation}
0<t<\min\left\{1,\frac{1-\lambda_0}{2k_*},
\frac{\Delta}{2(k_*M+k_b)}\right\}
\label{eq:gain_choice}
\end{equation}
works, omitting zero-denominator terms. Indeed, $\lambda\le\lambda_0+tk_F$, $\lambda_{\rm fr}\le\lambda_0+tk_{\rm fr}$ and $b_u\le b_0^{\rm drv}+tk_b$. The contraction increment uses less than half the original rate margin, while $t(k_*M+k_b)<\Delta/2$ leaves a strict radius margin. The ball is chosen before the gains. A residual bound available only on a smaller set cannot justify enlarging this ball. For arbitrary complete and frozen points in it, subtracting the top blocks gives the uniform mismatch bound $a_z(2g_AM_z+2g_BM_u)$ used in the analysis.

A nontrivial member of the mathematical class takes $n=m=1$, $P_h=2$, $\vz=(x,\psi)$, $\psi(X)=\alpha_{\rm d} x_{\rm past}$, and $\alpha_{\rm d}\ne0$. Set $c=a+v$ and
\begin{equation}
\mA(c)=\begin{bmatrix}a_0+g_A\tanh c&0\\\alpha_{\rm d}&0\end{bmatrix},\quad
\mB(x)=\begin{bmatrix}b_0+g_B\tanh x\\0\end{bmatrix},
\label{eq:analytic_example}
\end{equation}
where $|a_0|<1$, $b_0\ne0$, and both positive gains obey \eqref{eq:gain_choice}. These are rank-one corrections with genuine history and measurement dependence. Since $\psi^+=\alpha_{\rm d}x=\alpha_{\rm d}a^+$, the encoding graph is preserved. On an input box containing zero in its interior, a small alternating input $\vu_t^r=h(-1)^t$ has a unique periodic state in the already invariant ball: its two-step map is a contraction. At $h=0$, its linearization gives $x_t^r=-b_0h(-1)^t/(1+a_0)+o(h)$, so the measurement is nonconstant for sufficiently small nonzero $h$. This delay-coordinate example establishes nonemptiness of the model class, not exact representation or encoding consistency of an arbitrary Transformer.

Finally, the actual residual interface can be decomposed without assuming invariance of the encoder image. At an encoded initial point, let $\widehat\vz$ be the predicted next lift, $\widehat x=C_x\widehat\vz$, and $\widetilde X=\operatorname{shift}(\mX,\widehat x)$. If $C_\psi$ selects the learned feature block and $E_X$ appends a window-end perturbation, define
\begin{equation}
\begin{split}
\nu&=x^+-\widehat x,\qquad
\kappa=\psi_\theta(\widetilde X)-C_\psi\widehat\vz,\\
r&=\col\bigl(\nu,\kappa+
\psi_\theta(\widetilde X+E_X\nu)-\psi_\theta(\widetilde X)\bigr).
\end{split}
\label{eq:encoded_residual_decomposition}
\end{equation}
With exact caches, $w=\col(r,0,0)$. Uniform bounds $\bar\nu,\bar\kappa$ and a feature Lipschitz constant $L_\psi$ on all these windows give the sufficient budget $\bar w=a_z\sqrt{\bar\nu^2+(\bar\kappa+L_\psi\bar\nu)^2}$. The two terms measure physical measurement-prediction error and encoding-graph defect, respectively; neither is supplied by the gain construction alone.

\section{Technical Estimates for Frozen MPC}
\label{app:estimates}
For the condensed QP, stack $P=\col(p_{0|k},\ldots,p_{N|k})=P_0+\mathcal T U$. Let $K_pP=\col(E_{\rm prev}p_{0|k},\ldots,E_{\rm prev}p_{N-1|k})$, $\mathcal D_\Delta=I-K_p\mathcal T$, $Q_d=\operatorname{diag}(\Pi,\ldots,\Pi,\beta\Pi)$, $R_d=I_N\otimes R$, and $S_d=I_N\otimes S$. The Hessian of \eqref{eq:quadratic_template} is
\begin{equation}
\mathcal H=2(\mathcal T^TQ_d\mathcal T+R_d+
\mathcal D_\Delta^TS_d\mathcal D_\Delta)\succeq0.
\label{eq:QP_Hessian}
\end{equation}
Domain invariance follows from \eqref{eq:domain_radius} in the order established in the proof of Theorem~\ref{thm:tracking}: nominal rollout, frozen candidates, and then the next actual initialization. No estimate presupposes future membership.

For the strict convexity claim, let $\Delta U\ne0$ and let $j$ be its first nonzero block. Frozen causality gives $\Delta p_i=0$ for $i\le j$, while the last input-cache block of $\Delta p_{j+1}$ is $\Delta u_j\ne0$. If $j+1<N$, its state quadratic is weighted by $\Pi\succ0$; if $j+1=N$, it is weighted by $\beta\Pi\succ0$. Every other quadratic term is nonnegative. Thus the second variation is strictly positive. This argument neither assumes controllability of the lifted measurement block nor requires $R\succ0$.

Let $f_i$ and $p_i$ be full and frozen predictions for the same arbitrary candidate sequence. Both start at the point used for the first nominal coefficient evaluation, giving $f_0=p_0$ and $f_1=p_1$. Thereafter,
\begin{equation*}
\norm{f_{i+1}-p_{i+1}}_\Pi\le
\lambda\norm{f_i-p_i}_\Pi+\bar\eta,
\end{equation*}
by inserting $F(p_i,u_i)$. Hence
\begin{equation}
\begin{split}
\norm{f_i-p_i}_\Pi&\le D_i,\qquad D_0=D_1=0,\\
D_i&=\bar\eta\sum_{s=0}^{i-2}\lambda^s\quad(i\ge2).
\end{split}
\label{eq:frozen_discrepancy}
\end{equation}
For the nominal input sequence, induction gives identical full and frozen state sequences. The shifted comparison \eqref{eq:shift_comparison} uses the recurrence for $j<N-1$ and two complete updates at the reference tail, giving $\lambda d_{N-1}$.

The uniform shift bounds are
\begin{equation}
d_j=\lambda^j\bar w+\bar\eta\sum_{s=0}^{j-1}\lambda^s,\qquad 0\le j<N.
\label{eq:shift_error}
\end{equation}
With $\tau_{\rm ref}=\beta(1+1/\vartheta)\bar d_{\rm ref}^2$ and $L_\ell,L_f$ from \eqref{eq:cost_continuity}, the cost budget is
\begin{equation}
\varepsilon=\tau_{\rm ref}+L_\ell\sum_{j=0}^{N-1}d_j+L_f\lambda d_{N-1}.
\label{eq:epsilon_closed_loop}
\end{equation}
These constants are uniform. For an exact reference, use $\tau_{\rm ref}=0$ and the sharper terminal weight stated after Theorem~\ref{thm:tracking}.

To derive \eqref{eq:tail_decrease}, put $e=q-q_t^r$. At $u=\vu_t^r$, the $R$ term is zero and the weighted relative increment is $-S^{1/2}E_{\rm prev}e$, so $\ell_t\le c_S\norm{e}_\Pi^2$. By contraction and Young's inequality,
\begin{equation*}
\norm{F(q,\vu_t^r)-q_{t+1}^r}_\Pi^2
\le(1+\vartheta)\lambda^2\norm{e}_\Pi^2
+(1+1/\vartheta)\bar d_{\rm ref}^2.
\end{equation*}
The coefficient of $\norm{e}_\Pi^2$ after adding the stage cost and subtracting the old terminal cost is $c_S-\beta[1-(1+\vartheta)\lambda^2]\le0$. For $\bar d_{\rm ref}=0$ use $\lambda^2$ directly, without Young's inequality.

The explicit uniform cost bounds used in \eqref{eq:epsilon_closed_loop} are
\begin{equation}
\begin{split}
D_u^S&=\sup_{u,v\in\mathcal U}\norm{S^{1/2}(u-v)}_2,\quad
T_\Delta=D_u^S+2M\norm{H_S}_2,\\
L_\ell&=4M+2T_\Delta\norm{H_S}_2,\qquad L_f=4\beta M.
\end{split}
\label{eq:cost_continuity}
\end{equation}
For their derivation, errors relative to $q_t^r\in\mathcal Q$ have norm at most $2M$. Their squared-norm difference is at most $4M\norm{q-q'}_\Pi$. The weighted increment vector is
\begin{equation*}
S^{1/2}(u-\vu_t^r)-H_S\Pi^{1/2}(q-q_t^r),
\end{equation*}
whose norm is at most $T_\Delta$ and whose difference at a common input is at most $\norm{H_S}_2\norm{q-q'}_\Pi$. Thus its squared-norm difference is at most $2T_\Delta\norm{H_S}_2\norm{q-q'}_\Pi$. The $R$ term cancels, and the terminal estimate multiplies the state estimate by $\beta$, proving \eqref{eq:cost_continuity}.

\section{Value Comparison and Feasible Solver Error}
\label{app:comparison}
For $D_i$ in \eqref{eq:frozen_discrepancy}, define the reference/frozen error accumulation and the comparison constants by
\begin{equation}
\begin{split}
b_i&=\bar d_{\rm ref}\sum_{s=0}^{i-1}\lambda^s+D_i,\qquad 0\le i\le N,\\
\mathsf A&=2\left(\sum_{i=0}^{N-1}\lambda^{2i}+\beta\lambda^{2N}\right)
+\norm{H_S}_2^2,\\
\mathsf B&=2\left(\sum_{i=0}^{N-1}b_i^2+\beta b_N^2\right).
\end{split}
\label{eq:value_constants}
\end{equation}
The constants are uniform over all allowed nominal sequences; $\mathsf A\ge2$ gives $\rho=1-1/\mathsf A\in[0,1)$.

Consider the full trajectory for the reference input candidate from an arbitrary current state. Its reference error $r_i$ satisfies $r_{i+1}\le\lambda r_i+\bar d_{\rm ref}$, hence $r_i\le\lambda^i\norm{e_k}_\Pi+\bar d_{\rm ref}\sum_{s=0}^{i-1}\lambda^s$. Adding the full/frozen discrepancy $D_i$ gives the bound $\lambda^i\norm{e_k}_\Pi+b_i$. The reference candidate has zero input-error cost. For $i\ge1$ its predicted previous-input slot is exactly $\vu_{k+i-1}^r$, so the relative-increment cost is also zero; only the initial term, at most $\norm{H_S}_2^2\norm{e_k}_\Pi^2$, remains. Applying $(a+b)^2\le2a^2+2b^2$ to every state and terminal error yields \eqref{eq:value_comparison} with exactly \eqref{eq:value_constants}.

To expose the cost cancellation in \eqref{eq:value_descent}, write $g=F(p_N,\vu_{k+N}^r)$. The matched stages and the new terminal term satisfy
\begin{align*}
J_{k+1}^{\rm fr}(U^+;U^+)&\le\sum_{i=1}^{N-1}\ell_{k+i}(p_i,\widehat\vu_i)\\
&\quad+\ell_{k+N}(p_N,\vu_{k+N}^r)\\
&\quad+\beta\norm{g-q_{k+N+1}^r}_\Pi^2\\
&\quad+L_\ell\sum_{j=0}^{N-1}d_j+L_f\lambda d_{N-1}\\
&\le J_k^{\rm fr}(\widehat U_k;\bar U_k)
-\ell_k(p_0,\widehat\vu_0)+\varepsilon.
\end{align*}
The sum is empty for $N=1$. General cache perturbations are already covered by the state continuity estimates; no equality of all caches is used.

The comparison bound is for the exact value $W_k$ of each nominal-dependent QP, even when its delivered sequence is suboptimal. Combining the gap \eqref{eq:solver_gap} with the preceding inequality adds $\delta_k$, not $\delta_k+\delta_{k+1}$. Expanding the resulting scalar recursion gives
\begin{equation*}
W_k\le\rho^kW_0+\mathsf B(1-\rho^k)
+\sum_{t=0}^{k-1}\rho^{k-1-t}(\varepsilon+\delta_t).
\end{equation*}
Using $W_0\le\mathsf A\norm{e_0}_\Pi^2+\mathsf B$ gives $W_k\le\Xi_k$ with $\Xi_k$ defined in \eqref{eq:squared_tracking_bound}; the lower comparison then bounds the state error. In particular, the limiting input radius from \eqref{eq:input_tracking} is at most $\sqrt{[\mathsf B+\mathsf A(\varepsilon+\bar\delta)+\bar\delta]/\lambda_{\min}(R)}$ when $R\succ0$. A singular $R$ cannot supply this Euclidean input bound by division, although the state theorem and strict convexity remain valid.

A feasible output and a dual-feasible lower value $L_k\le W_k$ from the same QP give the valid gap bound $J_k^{\rm fr}(\widehat U_k)-L_k$; obtaining such a certificate does not require a second QP. More explicitly, for $J(U)=\tfrac12U^T\mathcal HU+g^TU+c$, constraints $G_UU\le h$, a certified $\mathcal H\succeq\mu I\succ0$, and $y\ge0$, let $r_{\rm stat}=\mathcal H\widehat U+g+G_U^Ty$. Strong convexity of the Lagrangian gives
\begin{equation*}
J(\widehat U)-W\le\frac{\norm{r_{\rm stat}}_2^2}{2\mu}
+y^T(h-G_U\widehat U).
\end{equation*}
Feasibility, $\mu$, and residual calculations must be valid under the numerical error model used. The theorem does not assert that a particular iteration limit or solver automatically meets these conditions, and an unanalysed infeasible fallback or input modification is outside the specified control law.

\begin{thebibliography}{99}
\bibitem{ljung1999system}
L. Ljung, \emph{System Identification: Theory for the User}, 2nd ed. Upper Saddle River, NJ, USA: Prentice Hall PTR, 1999.
\bibitem{nelles2020nonlinear}
O. Nelles, \emph{Nonlinear System Identification: From Classical Approaches to Neural Networks, Fuzzy Models, and Gaussian Processes}, 2nd ed. Cham, Switzerland: Springer, 2020.
\bibitem{forgione2021dynonet}
M. Forgione and D. Piga, ``dynoNet: A neural network architecture for learning dynamical systems,'' \emph{Int. J. Adapt. Control Signal Process.}, vol. 35, no. 4, pp. 612--626, 2021.
\bibitem{koopman1931hamiltonian}
B. O. Koopman, ``Hamiltonian systems and transformation in Hilbert space,'' \emph{Proc. Natl. Acad. Sci. USA}, vol. 17, no. 5, pp. 315--318, 1931.
\bibitem{mezic2005spectral}
I. Mezi\'c, ``Spectral properties of dynamical systems, model reduction and decompositions,'' \emph{Nonlinear Dyn.}, vol. 41, no. 1--3, pp. 309--325, 2005.
\bibitem{williams2015data}
M. O. Williams, I. G. Kevrekidis, and C. W. Rowley, ``A data-driven approximation of the Koopman operator: Extending dynamic mode decomposition,'' \emph{J. Nonlinear Sci.}, vol. 25, no. 6, pp. 1307--1346, 2015.
\bibitem{proctor2016dmdc}
J. L. Proctor, S. L. Brunton, and J. N. Kutz, ``Dynamic mode decomposition with control,'' \emph{SIAM J. Appl. Dyn. Syst.}, vol. 15, no. 1, pp. 142--161, 2016.
\bibitem{proctor2018generalizing}
J. L. Proctor, S. L. Brunton, and J. N. Kutz, ``Generalizing Koopman theory to allow for inputs and control,'' \emph{SIAM J. Appl. Dyn. Syst.}, vol. 17, no. 1, pp. 909--930, 2018.
\bibitem{korda2018linear}
M. Korda and I. Mezi\'c, ``Linear predictors for nonlinear dynamical systems: Koopman operator meets model predictive control,'' \emph{Automatica}, vol. 93, pp. 149--160, 2018.
\bibitem{brunton2022modern}
S. L. Brunton, M. Budi\v{s}i\'c, E. Kaiser, and J. N. Kutz, ``Modern Koopman theory for dynamical systems,'' \emph{SIAM Rev.}, vol. 64, no. 2, pp. 229--340, 2022.
\bibitem{lusch2018deep}
B. Lusch, J. N. Kutz, and S. L. Brunton, ``Deep learning for universal linear embeddings of nonlinear dynamics,'' \emph{Nat. Commun.}, vol. 9, no. 1, p. 4950, 2018.
\bibitem{takeishi2017koopman}
N. Takeishi, Y. Kawahara, and T. Yairi, ``Learning Koopman invariant subspaces for dynamic mode decomposition,'' in \emph{Proc. NeurIPS}, 2017, pp. 1130--1140.
\bibitem{otto2019linearly}
S. E. Otto and C. W. Rowley, ``Linearly-recurrent autoencoder networks for learning dynamics,'' \emph{SIAM J. Appl. Dyn. Syst.}, vol. 18, no. 1, pp. 558--593, 2019.
\bibitem{han2020deep}
Y. Han, W. Hao, and U. Vaidya, ``Deep learning of Koopman representation for control,'' in \emph{Proc. IEEE Conf. Decision Control (CDC)}, 2020, pp. 1890--1895.
\bibitem{kamb2020timedelay}
M. Kamb, E. Kaiser, S. L. Brunton, and J. N. Kutz, ``Time-delay observables for Koopman: Theory and applications,'' \emph{SIAM J. Appl. Dyn. Syst.}, vol. 19, no. 2, pp. 886--917, 2020.
\bibitem{das2019delay}
S. Das and D. Giannakis, ``Delay-coordinate maps and the spectra of Koopman operators,'' \emph{J. Stat. Phys.}, vol. 175, no. 6, pp. 1107--1145, 2019.
\bibitem{shenoy2025temporal}
S. Shenoy, B. Sharan, and H. Werner, ``Temporal autoencoder for identification and predictive control of nonlinear dynamics based on Koopman operator theory,'' in \emph{Proc. American Control Conf. (ACC)}, 2025, pp. 1795--1800.
\bibitem{han2024robust}
M. Han, Z. Li, Xiang Yin, and Xunyuan Yin, ``Robust learning and control of time-delay nonlinear systems with deep recurrent Koopman operators,'' \emph{IEEE Trans. Ind. Informat.}, vol. 20, no. 3, pp. 4675--4684, Mar. 2024.
\bibitem{forootani2026deepkoopformer}
A. Forootani, M. Khosravi, and M. Barati, ``DeepKoopFormer: a Koopman enhanced transformer based architecture for time series forecasting,'' \emph{npj Artificial Intelligence}, vol. 2, art. 35, 2026, doi: 10.1038/s44387-026-00085-3.
\bibitem{li2025mamko}
Z. Li, M. Han, and X. Yin, ``MamKO: Mamba-based Koopman operator for modeling and predictive control,'' in \emph{Proc. ICLR}, 2025.
\bibitem{liu2023koopa}
Y. Liu, C. Li, J. Wang, and M. Long, ``Koopa: Learning non-stationary time series dynamics with Koopman predictors,'' in \emph{Proc. NeurIPS}, 2023.
\bibitem{cisneros2016qlmpc}
P. S. G. Cisneros, S. Voss, and H. Werner, ``Efficient nonlinear model predictive control via quasi-LPV representation,'' in \emph{Proc. 55th IEEE Conf. Decision Control (CDC)}, 2016, pp. 3216--3221.
\bibitem{cisneros2020qlpv}
P. S. G. Cisneros and H. Werner, ``Nonlinear model predictive control for models in quasi-linear parameter varying form,'' \emph{Int. J. Robust Nonlinear Control}, vol. 30, no. 10, pp. 3945--3959, 2020.
\bibitem{diehl2005realtime}
M. Diehl, H. G. Bock, and J. P. Schl\"oder, ``A real-time iteration scheme for nonlinear optimization in optimal feedback control,'' \emph{SIAM J. Control Optim.}, vol. 43, no. 5, pp. 1714--1736, 2005.
\bibitem{gros2020rti}
S. Gros, M. Zanon, R. Quirynen, A. Bemporad, and M. Diehl, ``From linear to nonlinear MPC: Bridging the gap via the real-time iteration,'' \emph{Int. J. Control}, vol. 93, no. 1, pp. 62--80, 2020.
\bibitem{vaswani2017attention}
A. Vaswani \emph{et al.}, ``Attention is all you need,'' in \emph{Proc. NeurIPS}, 2017, pp. 5998--6008.
\end{thebibliography}
\end{document}
